\documentstyle[%


righttag%


,11pt%


,amssymb%


,russian%               remove in Eng.


,izvru%                 Set izven% in Eng.


]{amsart}





% {>}





\newcommand{\la}{\langle}


\newcommand{\ra}{\rangle}


\newcommand{\gf}{\frak f}


\newcommand{\gj}{\frak j}


\newcommand{\gh}{\frak h}


\newcommand{\loc}{\operatorname{loc}}


\newcommand\Mod{\operatorname{Mod}}


\newcommand\APM{\operatorname{APM}}


\newcommand\DIR{\operatorname{DIR}}


\newcommand\esssup{\operatorname{ess\,sup}}


\newcommand\diam{\operatorname{diam}}


\newcommand\frm{\operatorname{frm}}


\newcommand\rpm{\operatorname{rpm}}


\newcommand\comp{\operatorname{comp}}


\newcommand{\downto}{\downarrow}


\newcommand {\lj}{\lambda_k^{(j)}}


\newcommand{\tm}{(t,\omega)}


\newcommand{\mm}{\omega\in\Omega}


\newcommand{\tts}[1]{}





%\hoffset=-15mm%                  For Eng.


\setcounter{page}{34}


\god{2001}


\nomer{6~(469)} %                        Remove in Eng.


%\nomer{Vol.\,45, No.\,6}%               Set in Eng.


%\str{\pageref{begin}--\pageref{end}}%  Set in Eng.


\udk{517.934}





\author{\it А.Г.\,Иванов}


% NB:   ^^ remove in Eng.


\title{Об одном свойстве почти периодического


интеграла, зависящего от параметра}


\thanks{Работа выполнена при финансовой поддержке Российского фонда


фундаментальных исследований (гранты 97-01-00413 и 99-01-00454)


и Kонкурсного центра фундаментального


естествознания (грант\linebreak 97-0-1.9).}





\date{}


%\pagestyle{myheadings}%         Set in Eng.


\pagestyle{plain} %              Remove in Eng.


\begin{document}


%\thispagestyle{plain}%          Set in Eng.


\maketitle


\label{begin}





{\bf 1.} Обозначим через $B(\Bbb R,\frak Y)$ и


$S(\Bbb R,\frak Y)$, $\frak Y\subset


\Bbb R^n$, совокупность почти периодических (п.\,п.) по Бору,


соответственно по Степанову, функций $f:\Bbb R\to \frak Y$ [1].


Напомним, что для каждой п.\,п. (как по Бору, так и по


Степанову) функции $f$ существует среднее $M\{f(t)\}\doteq


\lim\limits_{T\to \infty}\frac{1}{T}\int\limits_0^T f(t)\,dt$


и, если $\Lambda (f)$ --- множество показателей Фурье п.\,п.


функции $f$, то $\Mod (f)\doteq \Mod (\Lambda (f))$


--- ее модуль, т.\,е. наименьшая группа по


сложению, содержащая множество $\Lambda (f)$. В дальнейшем


$(\Omega,\rho)$ --- компактное метрическое пространство и


$B(\Bbb R\times\Omega,\frak Y)$ --- совокупность отображений


\begin{equation}


\tm\mapsto f\tm\in\frak Y,


\quad\tm\in \Bbb R\times\Omega,


\tag{1.1}


\end{equation}


которые п.\,п. по $t$ в смысле Бора равномерно по $\omega\in \Omega$ [2].


Далее будем говорить, что отображение (1.1)


п.\,п. по $t\in\Bbb R$ в смысле Степанова равномерно по $\mm$


(пишем $f\in S(\Bbb R\times\Omega,\frak Y)$),


если оно удовлетворяет одновременно следующим условиям:


$f(\cdot,\omega)\in S(\Bbb R,\frak Y)$ при каждом $\omega$ и


$\lim\limits_{\gamma\downto 0}\frak d_{\gamma}[f,\Omega]=0$, где


\begin{equation}


\frak d_{\gamma}[f,\Omega]\doteq\sup


\begin{Sb}


\omega_1,\omega_2 \in\Omega\\ \rho(\omega_1,\omega_2)\le \gamma\end{Sb}


d(f(\cdot,\omega_1),f(\cdot,\omega_2))\tag{1.2}


\end{equation}


(определение метрики $d$ см. в [1]).


Если $f\in B(\Bbb R\times\Omega,\frak Y)$


(или $f\in S(\Bbb R\times\Omega,\frak Y)$),


то множество $\Lambda(f)\doteq


\bigcup\limits_{\omega\in\Omega}\Lambda (f(\cdot,\omega))$ назовем


множеством показателей Фурье отображения $f$.





Введем теперь мерозначные п.\,п. функции. С этой целью


обозначим через $(\frm (U),|\cdot |_{w})$ (см. [3]) нормированное


пространство таких мер Радона на $\Bbb R^m$, носитель которых


содержится во множестве $U\in \comp (\Bbb R^m)$, и через $\rpm (U)$


--- его подмножество, состоящее из вероятностных мер Радона. В


дальнейшем $\DIR (U)$ --- совокупность мер Дирака $\delta_u$,


сосредоточенных в точках $u\in U$, и $\cal M=\cal M(\Bbb R,\frm(U))$


--- множество таких измеримых отображений $\mu: \Bbb R\to (\frm


(U), |\cdot |_w)$, что $\|\mu\|\doteq\esssup\limits_{t\in \Bbb R}


|\mu(t)|(U)<\infty$ ($|\mu(t)|(U)$ --- вариация меры $\mu(t)$).


Пусть далее $\frak B_n=\frak B(\Bbb R\times U,\Bbb R^n)$ --- совокупность


отображений $\varphi:\Bbb R\times U\to


\Bbb R^n$, удовлетворяющих следующим условиям:


при почти всех (п.\,в.) $t\in\Bbb R$\;


$\varphi (t,\cdot)\in C(U,\Bbb R^n)$, для каждого $u\in U$


отображение $t\mapsto \varphi(t,u)\in \Bbb R^n$, $t\in \Bbb R$,


измеримо и существует такая функция $\psi_{\varphi}\in L_1


(\Bbb R,\Bbb R_+)$, что при п.\,в. $t\in \Bbb R$\;


$\max\limits_{u\in U}|\varphi (t,u)|\le \psi_{\varphi}(t)$.


Несложно показать, что отображение $\varphi\mapsto


\|\varphi\|_{\frak B_n}\doteq \int\limits_{\Bbb R}\max\limits_{u\in


U}|\varphi(t,u)|dt$, $\varphi \in \frak B_n$,


является нормой в $\frak B_n$, и


незначительным изменением схемы доказательства теоремы


Данфорда--Петтиса ([3], c.\,299) можно показать, что


$\cal M\cong\frak B_1^*$. Это позволяет [4]--[6] ввести в $\cal M$


норму $\|\cdot\|_w$, относительно которой множества


$\cal M_1\doteq\cal M(\Bbb R,\rpm(U))$ и $\Sigma_1\doteq


\{\mu\in\cal M : \|\mu\|\le 1\}$ компактны.


\medskip





{\bf Определение 1.1} [4]. Отображение $\mu\in\cal M$


называется п.\,п. по Степанову, если для любой функции $c\in


C(U,\Bbb R)$ отображение


$$


t\mapsto\la \mu(t),c(u)\ra\doteq \int_U c(u)\mu(t)(du)


$$


принадлежит пространству $S(\Bbb R,\Bbb R)$.


\medskip





Совокупность всех п.\,п. по Степанову отображений из


$\cal M$ обозначим $\APM$, и $\APM_1\doteq\APM\cap\cal M_1$. Далее,


через $\APM_1^{(1)}$ обозначим совокупность $\mu\in


\APM_1$ таких, что $\mu(t)=\delta_{u(t)}$ при п.\,в. $t\in \Bbb R$ и


некотором измеримом отображении $u:\Bbb R\to U$. Можно


показать, что $S(\Bbb R,U)\cong \APM_1^{(1)}$, и, следовательно, каждое


$u(\cdot)\in S(\Bbb R,U)$ можно рассматривать как элемент множества


$\APM_1^{(1)}\subset \APM$,


отождествляя его с отображением $t\to


\delta_{u(t)}\in \DIR(U)$. Каждому $\mu\in \APM$ можно


однозначно поставить в соответствие мерозначный ряд Фурье


[4]--[6] и определить множество $\Lambda(\mu)$ его


показателей Фурье, и $\Mod


(\mu)\doteq \Mod(\Lambda(\mu))$, если $\mu\in \APM$.


\medskip





{\bf Определение 1.2.} Отображение $\tm\to


\mu\tm\in \frm(U)$ называется п.\,п. по $t\in \Bbb R$ в смысле Степанова


равномерно по $\mm$ (пишем $\mu\in S(\Bbb R\times\Omega,\frm (U))$),


если для каждой функции $c\in C(U,\Bbb R)$ отображение


$$


\tm\to \la \mu(t,\omega),c(u)\ra\doteq


\int_{U}c(u)\mu\tm(du)


$$


принадлежит пространству $S(\Bbb R\times\Omega,\Bbb R)$.


\medskip





Из определений 1.1 и 1.2 вытекает, что 


$\mu(\cdot,\omega)\in \APM$ при каждом $\mm$. Поэтому следующее множество


$\Lambda(\mu)=\bigcup\limits_{\omega\in\Omega}\Lambda(\mu(\cdot,\omega))$


естественно назвать множеством показателей Фурье отображения


$\mu\in S(\Bbb R\times\Omega, \rpm(U))$,


а множество $\Mod(\mu)\doteq \Mod(\Lambda(\mu))$ --- его модулем.





В дальнейшем каждую функцию $g$ из $L_{1}^{\loc}(\Bbb R, C(U,\Bbb R))$


представляем в виде отображения $(t,u)\mapsto g(t,u)\in


\Bbb R,(t,u)\in\Bbb R\times U$,


и через $S(\Bbb R,C(U,\Bbb R))$ обозначим совокупность


таких функций $g$ из $L_{1}^{\loc}(\Bbb R,C(U,\Bbb R))$,


что для любого $\varepsilon>0$ множество


$$


E_{S}(f,\varepsilon)\doteq


\bigg\{\tau\in\Bbb R:\sup\limits_{t\in\Bbb R}


\int_{t}^{t+1}\max\limits_{u\in U}


|g(s+\tau,u)-g(s,u)|ds\le\varepsilon\bigg\}


$$


относительно плотно.


Отметим, что $S(\Bbb R,C(U,\Bbb R))\subset S(\Bbb R\times U,\Bbb R)$.


\medskip





{\bf Лемма 1.1.} {\it


Если $g\in S(\Bbb R,C(U,\Bbb R))$, то


$\lim\limits_{\gamma\downarrow 0}


\Big(\sup\limits_{t\in\Bbb R}\int\limits_{t}^{t+1}


w_{\gamma}[g(s,\cdot),U] ds\Big)=0$,


где $w_{\gamma}[g(s,\cdot),U]$ --- $\gamma$-колебание


функции $u\mapsto g(s,u)$ на множестве $U$.


}





\begin{pf}


Так как $g\in S(\Bbb R,C(U,\Bbb R))$,


то для заданного $\varepsilon>0$ множество


$E_{S}(f,\varepsilon/3)$ относительно плотно


и, следовательно, найдется такое $l>0$,


что при каждом $t\in\Bbb R$ существует


$\tau\in [-t,-t+l]\cap E_{S}(g,\varepsilon/3)$.


Поэтому при каждом $t\in\Bbb R$


$$


\int_{t}^{t+1}w_{\gamma}[g(s,\cdot),U] ds\le


2\sup\limits_{t\in\Bbb R}


\int_{t}^{t+1}\max\limits_{u\in U}|g(s+\tau,u)-g(s,u)|ds


+\gf(\gamma)<2\varepsilon/3+\gf(\gamma)<2,


$$


где $\gf(\gamma)\doteq\int\limits_{0}^{l+1}w_{\gamma}[g(s,\cdot),U] ds$.


Используя


теорему, приведенную в ([3], с.\,158), несложно показать, что


$\lim\limits_{\gamma\downarrow 0}\gf(\gamma)=0$.


Отсюда совместно с приведенными выше соотношениями получаем нужное


равенство.


\end{pf}





{\bf Теорема 1.1.} {\it


Пусть $(\Omega,\rho)$ --- компактное метрическое


пространство, множество


$\frak A\doteq\{\mu(\cdot,\omega), \omega\in\Omega\}\subset


APM_{1}$ и равностепенно п.\,п.\footnote{\,т.\,е. для каждой функции


$c\in C(U,\Bbb R)$ совокупность отображений $\{\la\mu(\cdot,\omega),c(u)\ra,


\omega\in\Omega\}\subset S(\Bbb R,\Bbb R)$ равностепенно п.\,п.}.


Тогда для любой функции


$g\in S(\Bbb R,C(U,\Bbb R))$ совокупность отображений $\{t\mapsto


f(t,\omega)\doteq\la\mu(t,\omega),g(t,u)\ra, \omega\in\Omega\}$, где


$$


\la\mu(t,\omega),g(t,u)\ra\doteq\int_{U}g(t,u)\mu(t,\omega)(du),


$$


принадлежит $S(\Bbb R,\Bbb R)$


и является равностепенно п.\,п. Кроме того, если


$\mu\in S(\Bbb R\times\Omega,\rpm (U))$,


то функция $(t,\omega)\mapsto f(t,\omega)$


принадлежит пространству


$S(\Bbb R\times\Omega,\Bbb R)$, и еe модуль содержится


в $\Mod(\Lambda(\mu)\cup\Lambda(g))$.


}





\begin{pf}


Первое утверждение теоремы


1.1 можно доказать, следуя схеме доказательства теоремы 3.1 из [7]


(см. также [5]). Докажем второе


утверждение. Так как $g\in S(\Bbb R,C(U,\Bbb R))$, то по лемме 1.1 для


заданного $\varepsilon{>}0$ найдется такое $\gamma{>}0$, что


$\sup\limits_{t\in\Bbb R}\int\limits_{t}^{t+1}


w_{\gamma}[g(s,\cdot),U]ds<\varepsilon/6$.


Пусть, далее,


$\cal U_1,\dotsc,\cal U_{p}$ --- такое открытое покрытие компакта $U$,


что $\max\limits_{1\le j\le p} (\diam \cal U_{j})\le\gamma$, и через


$\{\alpha_{j}\}_{j=1}^{p}$ обозначим непрерывное


разбиение единицы, подчиненное этому покрытию.


Теперь для каждого $j=1,\dotsc,p$ фиксируем точку


$u_{j}\in U\cap\cal U_{j}$,


в которой $\alpha_{j}(u_{j})>0$, и рассмотрим отображение


$(t,\omega)\mapsto\lambda_{j}(t,\omega)\doteq


\la\mu(t,\omega),\alpha_{j}(u)\ra \in [0,1]$,


$(t,\omega)\in\Bbb R\times\Omega$.


Так как $\mu\in S(\Bbb R\times\Omega,\rpm(U))$,


то (см. определение 1.2) 


$\lambda_{j}\in S(\Bbb R\times\Omega,\Bbb R)$


при каждом $j=1,\dotsc,p$ и, кроме того,


$\sum\limits_{j=1}^{p}\lambda_{j}(t,\omega)=1$,


$(t,\omega)\in\Bbb R\times\Omega.$


Далее, для каждой функции


$g(\cdot,u_{j})\in S(\Bbb R,\Bbb R)$, $j=1,\dotsc,p$,


возьмем такую функцию


$\gf_{j}\in S(\Bbb R,\Bbb R)$ (см. [1], с.\,231), что


$\esssup\limits_{t\in \Bbb R}|\gf_{j}(t)|\doteq\frak k_{j}<\infty$ и


$d(g(\cdot,u_{j}),\gf_{j}(\cdot))<\varepsilon/3p$. Поскольку


$\lambda_{j}\in S(\Bbb R\times\Omega,\Bbb R)$, $j=1,\dotsc,p$, то


найдется такое $\wh\gamma\in(0,\gamma)$,


что при всех $\beta\in(0,\wh\gamma)$


будет выполнено неравенство (см. (1.2))


$\frak d_{\beta}[\lambda_{j},\Omega]<\varepsilon/3p\frak k$,


где $\frak k\doteq \max\limits_{1\le j\le p}\frak k_{j}$.


Тогда при этих $\beta$ получаем


$$


\frak d_{\beta}[f,\Omega]\le 2 \sup\limits_{t\in\Bbb R}\int_{t}^{t+1}


w_{\gamma}[g(s,\cdot),U]ds+\sum_{j=1}^{p}d(g(\cdot,u_{j}),\gf_{j}(\cdot))+


\frak k\sum_{j=1}^{p}\frak d_{\beta}[\lambda_{j},\Omega]<\varepsilon.


$$


Отсюда и из утверждения первой части теоремы 1.1 (см.


определение 1.2) вытекает, что


$f\in S(\Bbb R\times\Omega,\Bbb R)$. Доказательство, что


модуль функции $f$ содержится в $\Mod(\Lambda(\mu)\cup\Lambda(g))$,


аналогично доказательству соответствующего утверждения в теореме 3.1


из [7].


\end{pf}





{\bf Лемма 1.2.} {\it


Функция $u\in S(\Bbb R\times\Omega,U)$ в том и только том случае,


если отображение $\tm\mapsto \delta_{u\tm}$ принадлежит


$S(\Bbb R\times\Omega,\rpm(U))$ и их модули совпадают.


}


\medskip





Доказательство вытекает из определения пространства


$S(\Bbb R\times\Omega,U)$ и неравенства (см. (1.2))


$$


\sup\begin{Sb}


\omega_1,\omega_2\in \Omega \\ \rho(\omega_1,\omega_2)\le \gamma


\end{Sb}


\bigg(\sup\limits_{t\in\Bbb R}


\int_t^{t+1}|c(u(s,\omega_1))-c(u(s,\omega_2))|ds\bigg)\le


\frac{2}{\sigma}\|c\|_{C(U,\Bbb R)}\cdot


\frak d_{\gamma}[u,\Omega]+w_{\sigma}[c,U],


$$


справедливого для каждой


функции $c\in C(U,\Bbb R)$ и фиксированных констант $\sigma,\gamma>0$, где


$w_{\sigma}[c,U]$ --- $\sigma$-колебание функции $c(\cdot)$ на


множестве $U$.





Из леммы 1.2 и теоремы 1.1 вытекает


\medskip





{\bf Следствие 1.1.} Пусть $g\in S(\Bbb R,C(U,\Bbb R))$ и


$u\in S(\Bbb R\times\Omega,U)$.


Тогда отображение $\tm\mapsto g(t,u\tm)$ принадлежит


$S(\Bbb R\times\Omega,\Bbb R)$, и его модуль


содержится в $\Mod(\Lambda(g)\cup\Lambda(u))$.


\medskip





{\bf Следствие 1.2.} Пусть отображение $g:\Bbb R\times U\to \Bbb R$


принадлежит либо пространству


$S(\Bbb R,C(U,\Bbb R))$, либо $B(\Bbb R\times U,\Bbb R)$.


Тогда для всякого $\mu\in \APM_1$ и $u\in S(\Bbb R,U)$ функции $t\mapsto\la


\mu(t),g(t,u)\ra$, $t\mapsto g(t,u(t))$ принадлежат пространству


$S(\Bbb R,\Bbb R)$, и их модули содержатся в $\Mod


(\Lambda(\mu)\cup\Lambda(g))$ и $\Mod (\Lambda(g)\cup\Lambda(u))$


соответственно.


\medskip





Целью данной работы является


\medskip





{\bf Теорема 1.2.} {\it


Пусть $(\Omega,\rho)$ --- компактное метрическое


пространство, функция $c\in C(U,\Bbb R^n)$ и отображение $\mu\in


S(\Bbb R\times\Omega,\rpm(U))$ такое, что при каждом $\mm$ уравнение


\begin{equation}


\Dot x=\la


\mu(t,\omega),c(u)\ra,\ \ t\in \Bbb R,\tag{1.3}


\end{equation}


имеет п.\,п. по Бору решение


$x\tm=\int\limits_0^t\la\mu(s,\omega),c(u)\ra ds$, $t\in \Bbb R$.


Тогда существует такая последовательность функций


$\{v_i\}_{i=1}^{\infty}\subset S(\Bbb R\times\Omega,U)$, что


${\Mod}(v_i)\subset {\Mod}(\mu)$


при каждом $i\in \Bbb N$, и уравнение $\Dot x=c(v_i\tm)$,


$t\in \Bbb R$, $\mm$, имеет п.\,п. по Бору решение


$x_i\tm=\int\limits_0^t


c(v_i(s,\omega))ds$, $t\in\Bbb R$. При этом множество функций


$\{x_{i}(\cdot,\omega), i\in\Bbb N, \omega\in\Omega\}


\subset B(\Bbb R,\Bbb R^{n})$


ограничено и


}


$$


\lim\limits_{i\to\infty}(\sup_{\mm} M\{|x\tm-x_i\tm|\})=0.


$$





Доказательству теоремы 1.2 предпошлем ряд вспомогательных


утверждений.


\medskip





{\bf 2.} Имеет место


\medskip





{\bf Теорема 2.1.} {\it


Пусть $(\Omega,\rho)$ --- компактное


метрическое пространство. Тогда для каждого отображения $\mu\in


S(\Bbb R\times\Omega,\rpm(U))$ существует последовательность функций


$\{u_j\}_{j=1}^{\infty}$ из пространства $S(\Bbb R\times\Omega,U)$,


для которой


$\Mod(u_j)\subset\Mod(\mu)$ при всех $j\in\Bbb N$


и которая обладает следующими


свойствами\/$:$ $1)$~имеет место равенство


$\lim\limits_{j\to\infty}(\sup\limits_{\mm}\|\mu(\cdot,\omega)-


\delta_{u_j(\cdot,\omega)}\|_w)=0;$


$2)$~при каждом $j\in\Bbb N$


$$


\lim\limits_{\gamma\downto 0}\bigg(\sup\begin{Sb}


\omega_1,\omega_2\in\Omega \\ \rho(\omega_1,\omega_2)\le\gamma \end{Sb}


\bigg(\sup\limits_{t\in\Bbb R}


\int_t^{t+1}|\delta_{u_j(s,\omega_1)}-\delta_{u_j(s,\omega_2)}|


(U)\bigg)\bigg)ds=0;


$$


$3)$ для всякой функции $g\in S(\Bbb R,C(U,\Bbb R))$


\begin{gather*}


\sup\limits_{t\in\Bbb R}\bigg|\int_t^{t+1}


\la\mu(s,\omega)-\delta_{u_j(s,\omega)},g(s,u)\ra ds\bigg|


\underset{\mm}{\rightrightarrows} 0


\text{\; при \;}j\to\infty, \\


%


M\{g(t,u_j\tm)\}


\underset{\mm}{\rightrightarrows}


M\{\la \mu\tm,g(t,u)\ra\}


\text{\; при \;} j\to\infty.


\end{gather*}


}





Доказательство теоремы 2.1, а также доказательство нижеприводимой леммы 2.1,


аналогичны доказательствам соответствующих утверждений,


приведенных в [5], для отображения


$(t,\omega)\mapsto\mu(t,\omega)\in\rpm(U)$ такого, что (ср. с определением


1.2) для каждой функции $c\in C(U,\Bbb R)$ отображение


$(t,\omega)\mapsto\la\mu(t,\omega),c(u)\ra$ принадлежит пространству


$S(\Bbb R,C(U,\Bbb R))$.


Поэтому приведем лишь кратко необходимую в


дальнейшем конструкцию функций $u_{j}(\cdot)$ из теоремы~2.1.





Для каждого $j\in \Bbb N$ строим такое открытое покрытие


$\cal U_1^{(j)},\dotsc,\cal U_{p_j}^{(j)}$ компакта $U$,


что\linebreak $\max\limits_{1\le k\le p_j} (\diam


\cal U_k^{(j)})\le 1/j$, и через


$\{\alpha_k^{(j)}\}_{k=1}^{p_j}$ обозначим непрерывное


разбиение единицы, подчиненное этому покрытию.


Теперь для каждого $k=1,\dotsc,p_j$ фиксируем точку


$u_k^{(j)}\in U\cap\cal U_k^{(j)}$,


в которой $\alpha_k^{(j)}(u_k^{(j)})>0$,


и рассмотрим функцию


${(t,\omega)\mapsto\lambda_k^{(j)}\tm\doteq


\la \mu\tm,\alpha_k^{(j)}(u)\ra \in [0,1]}$, ${\tm\in \Bbb R\times\Omega}$.


Так как $\mu\in S(\Bbb R\times\Omega,\rpm(U))$, то (см. определение 1.2)


$\{\lj\}_{k=1}^{p_j}\subset S(\Bbb R\times\Omega,[0,1])$ и при этом


$\sum\limits_{k=1}^{p_j}\lj\tm=1$,


$\tm\in \Bbb R\times\Omega$. При каждом


$j\in\Bbb N$ введем отображение $\tm\to\Delta_j \tm\doteq


\sum\limits_{k=1}^{p_j}\lj\tm\delta_{u_k^{(j)}}\in \rpm(U)$,


$\tm\in\Bbb R\times\Omega$.


Можно показать, что при каждом


$j\in\Bbb N$\; $\Delta_j\in S(\Bbb R\times\Omega,\rpm(U))$ и


$\Mod(\Delta_j)\subset\Mod(\mu)$. Выбираем далее число $a>0$ таким, чтобы


$4\pi/a\in\Mod(\mu)$ и отрезок $[0,a]$ разбиваем на $j$ равных


отрезков $I_l^{(j)}=\big[\frac{l-1}{j}a,\frac{l}{j}a\big]$, $l=1,\dotsc,j$.


В свою очередь, каждый отрезок $I_l^{(j)}$ разобьем на $p_j$


частичных подотрезков $I_{l_k}^{(j)}(\xi,\omega)$,


$k=1,\dotsc,p_j$, зависящих от


$(\xi,\omega)\in \Bbb R\times\Omega$,


определенных равенствами


\begin{equation}


\begin{split}


I_{l_1}^{(j)}(\xi,\omega) &\doteq


\frac{l-1}{j}a+\bigg[0,\int_{I_l^{(j)}}\lambda_1^{(j)}(t+\xi,


\omega)dt\bigg],\\


%


I_{l_k}^{(j)}(\xi,\omega) &\doteq


\frac{l-1}{j}a+\bigg[\sum\limits_{p=1}^{k-1}


\int_{I_l^{(j)}}\lambda_p^{(j)}(t+\xi,\omega)dt,\


\sum\limits_{p=1}^{k}


\int_{I_l^{(j)}}\lambda_p^{(j)}(t+\xi,\omega)dt\bigg],\quad k=2,\dotsc,p_j.


\end{split}


\tag{2.1}


\end{equation}


Из (2.1) следует, что отрезки


$I_{l_1}^{(j)}(\xi,\omega),\dotsc,I_{l_{p_j}}^{(j)}(\xi,\omega)$ примыкают


друг к другу и при каждом $l=1,\dotsc,j$ отрезок $I_l^{(j)}$ совпадает


с их объединением. Теперь рассмотрим последовательность


$\{w_m^{(j)}\}_{m\in\Bbb Z}$, состоящую


из отображений $w_m^{(j)}:[0,a]\times\Omega\to U$, $m\in\Bbb Z$,


определенных равенством


\begin{equation}


w_m^{(j)}\tm\doteq


\sum\limits_{l=1}^j


\chi_{I_l^{(j)}}(t)\sum\limits_{k=1}^{p_j}\chi_{I_{l_k}^{(j)}(ma,\omega)}(t)


u_k^{(j)},\quad \tm\in [0,a]\times \Omega,\tag{2.2}


\end{equation}


где $I_{l_k}^{(j)}(ma,\omega)$,


$k=1,\dotsc,p_j$, задаются равенствами (2.1) при $\xi=ma$, а


$\chi_A (\cdot)$ --- характеристическая функция множества


$A\subset \Bbb R$.


\medskip





{\bf Определение 2.1.} Последовательность


$\{\gf_m\}_{m\in\Bbb Z}$ отображений


\begin{equation}


\tm\to\gf_m\tm\in\frak Y,\quad \tm\in


[0,a]\times\Omega, \tag{2.3}


\end{equation}


называется п.\,п. равномерно по $\mm$, если


при каждом $\mm$ последовательность


$\{\gf_m(\cdot,\omega)\}_{m\in\Bbb Z}$


содержится в $L_1([0,a],\frak Y)$, является п.\,п. (т.\,е. для любого


$\varepsilon>0$ множество $\cal E


(\{\gf_m(\cdot,\omega)\}_{m\in\Bbb Z},\varepsilon)\doteq


\Big\{n\in\Bbb Z:


\sup\limits_{m\in\Bbb Z}


\int\limits_0^a |\gf_{m+n}\tm-\gf_m\tm|dt<\varepsilon\Big\}$


относительно плотно), и, кроме того, $\lim\limits_{\gamma\downto


0}\frak d_{\gamma}[\{\gf_m\}_{m\in\Bbb Z},\Omega]=0$, где


$$


\frak d_{\gamma}[\{\gf_m\}_{m\in\Bbb Z},\Omega]\doteq


\sup\begin{Sb}


\omega_1,\omega_2\in\Omega \\ \rho(\omega_1,\omega_2)\le \gamma\end{Sb}


\bigg(\sup\limits_{m\in\Bbb Z}


\int_0^a |\gf_{m}(t,\omega_1)-\gf_m(t,\omega_2)|dt\bigg).


$$


Если последовательность $\{\gf_m\}_{m\in\Bbb Z}$


отображений (2.3) п.\,п. равномерно по


$\mm$, то согласно определению 2.1 при каждом $\mm$ последовательность


$\{\gf_m(\cdot,\omega)\}_{m\in\Bbb Z}$ из $L_1([0,a],\frak Y)$ п.\,п.,


и, следовательно,


для нее можно определить множество показателей Фурье последовательности


$\Lambda(\{\gf_m(\cdot,\omega)\}_{m\in\Bbb Z})$. Тогда множество


$\Lambda(\{\gf_m\}_{m\in\Bbb Z})\doteq


\bigcup\limits_{\mm}\Lambda(\{\gf_m(\cdot,\omega)\}_{m\in\Bbb Z})$


называется множеством показателей Фурье последовательности


$\{\gf_m\}_{m\in\Bbb Z}$,


а


$\Mod(\{\gf_m\}_{m\in\Bbb Z})\doteq \Mod(\Lambda(\{\gf_m\}_{m\in\Bbb Z}))$


--- ее модулем.





Так же, как и в [5], показываем, что последовательность


$\{w_m^{(j)}\}_{m\in\Bbb Z}$


отображений (2.2) является п.\,п. равномерно по $\mm$ и


\begin{equation}


\Mod(\{w_m^{(j)}\}_{m\in\Bbb Z}) \subset a\Mod (\mu)+2\pi\Bbb Z.\tag{2.4}


\end{equation}


Отсюда (см. определения 1.2 и 2.1) несложно заключить, что


при каждом $j\in\Bbb N$ функция $u_j:\Bbb R\times\Omega\to U$,


определенная на каждом множестве $[ma,(m+1)a]\times\Omega$, $m\in\Bbb Z$,


равенством


$u_j(t+ma,\omega)\doteq w_m^{(j)}\tm$, $\tm\in [0,a]\times\Omega$,


принадлежит $S(\Bbb R\times\Omega,U)$,


и, кроме того, в силу выбора числа $a>0$ и


включения (2.4) $\Mod(u_j)\subset\Mod(\mu)$.





Полагаем


$$


\eta_j^{(1)}\tm\doteq \mu\tm-\Delta_j\tm,


\quad \eta_j^{(2)}\tm\doteq


\Delta_j\tm-\delta_{u_j\tm},\quad \tm\in\Bbb R\times\Omega,\ \; j\in\Bbb N.


$$


Ясно, что для всех $j\in \Bbb N$\;


$\eta_j^{(\frak q)}\in S(\Bbb R\times\Omega,\frm(U)),


\|\eta_j^{(\frak q)}\|\le


2$, $\frak q=1,2$. Так как ${\mu\in S(\Bbb R\times\Omega,\rpm(U))}$ и


$\diam\cal U_k^{(j)}\le 1/j$, то из соотношений


$$


\sup\limits_{t\in


\Bbb R}\int_t^{t+\frak l}|\la \eta_j^{(1)}(s,\omega),c(u)\ra|ds\le


\sup\limits_{t\in \Bbb R}\int_t^{t+\frak l}\bigg(\sum_{k=1}^{p_j}


\int_U\alpha_k^{(j)}(u)|c(u_k^{(j)})-c(u)|\mu(s,\omega)(du)\bigg)ds\le


\frak l w_{\frac{1}{j}}[c,U]


$$


вытекает, что при каждом $\frak l>0$


\begin{equation}


\sup\limits_{t\in \Bbb R}


\int_t^{t+\frak l}|\la\eta_j^{(1)}(s,\omega),c(u)\ra|ds


\underset{\mm}{\rightrightarrows} 0


\text{\; при \;}j\to\infty. \tag{2.5}


\end{equation}





{\bf Лемма 2.1.} {\it


Для каждой функции $g\in S(\Bbb R,C(U,\Bbb R))$


и любого фиксированного $\frak l>0$


\begin{equation}


\sup\limits_{t\in \Bbb R}\bigg|\int_t^{t+\frak l}\la


\eta_j^{(2)}(s,\omega),g(s,u)\ra ds\bigg|


\underset{\mm}{\rightrightarrows} 0


\text{\; при \;}j\to\infty. \tag{2.6}


\end{equation}


}


\medskip





{\bf Следствие 2.1.} Пусть $c\in C(U,\Bbb R)$. Тогда при любом


фиксированном $\frak l>0$


\begin{equation}


\sup\limits_{m\in\Bbb Z}\bigg(\sup\limits_{t\in[m\frak l,(m+1)\frak l]}


\bigg|\int_{m\frak l}^t\la\eta_j^{(2)}(s,\omega),c(u)\ra ds\bigg|\bigg)


\underset{\mm}{\rightrightarrows} 0


\text{\; при \;}j\to\infty. \tag{2.7}


\end{equation}


\medskip





{\bf Доказательство.}


Допустим противное. Тогда найдутся константа $\gamma>0$,


строго возрастающая последовательность


$\{j(N)\}_{N=1}^{\infty}\subset\Bbb N$, а


также последовательности ${\{\omega_N\}_{N=1}^{\infty}\subset\Omega,}$


$\{m_N\}_{N=1}^{\infty}\subset\Bbb Z$ и $\{t_N\}_{N=1}^{\infty}$,


где $t_N\doteq m_N\frak l+\theta_N\frak l$,


$\theta_N\in[0,1)$, такие, что для каждого $N\in\Bbb N$


будет выполнено неравенство


$$


\bigg|\int_{m_N\frak l}^{t_N}\la\eta_{j(N)}^{(2)}(s,\omega_N),


c(u)\ra ds\bigg|\ge \gamma.


$$


Поскольку $\{\theta_N\}_{N=1}^{\infty}\subset[0,1)$, то можно считать, что


$\lim\limits_{N\to\infty}\theta_N=\theta\in[0,1]$.


Поэтому, учитывая $|\la\eta_{j}^{(2)}(t,\omega),c(u)\ra|\le


2\|c\|_{C(U,\Bbb R)}$,


$\tm\in\Bbb R\times\Omega$, $j\in \Bbb N$, из указанного выше


неравенства получим существование $\wh N\in\Bbb N$, начиная с которого


будут выполняться неравенства


$$


\sup\limits_{\mm}\bigg(\sup\limits_{m\in\Bbb Z}


\bigg|\int_{m\frak l}^{m\frak l+\frak l}


\la\eta_{j(N)}^{(2)}(s,\omega),g(s,u)\ra ds\bigg|\bigg)\ge \gamma/2,


$$


где $g(t,u)\doteq\gf(t) c(u)$, $(t,u)\in \Bbb R\times U$,


а $\gf:\Bbb R\to\Bbb R$ --- $\frak l$-периодическая функция,


которая на $[0,\frak l]$ определяется равенством


$\gf(t)\doteq \chi_{[0,\theta\frak l]}(t)$, $t\in[0,\frak l]$.


Очевидно, так заданная функция $g$ принадлежит $S(\Bbb R,C(U,\Bbb R))$


и, следовательно, полученное неравенство несовместно с предельным


соотношением (2.6).





Далее, при каждом $\mm$ и любых $m\in\Bbb Z$,


$j\in\Bbb N$ (см. определение мер $\Delta_j$


и (2.2)) имеем


\begin{multline*}


\int_{ma}^{(m+1)a}\la \eta_j^{(2)}(s,\omega),c(u)\ra ds=


\sum_{l=1}^j\bigg(\int_{I_l^{(j)}}\sum\limits_{k=1}^{p_j}\lambda_k^{(j)}


(t+ma,\omega)


c(u_k^{(j)})dt-\sum\limits_{k=1}^{p_j}\int_{I_{l_k}^{(j)}(ma,\omega)}


c(u_k^{(j)})dt\bigg)= \\


%


=\sum\limits_{l=1}^j\sum\limits_{k=1}^{p_j}\bigg(\int_{I_l^{(j)}}


\lambda_k^{(j)}(t+ma,\omega)dt-\text{mes}\,


I_{l_k}^{(j)}(ma,\omega)\bigg)c(u_k^{(j)})\overset{(2.1)}{=}0,


\end{multline*}


поэтому при каждом $\mm$, $j\in \Bbb N$, $q\in \Bbb Z$


\begin{equation}


\int_0^{qa}\la\Delta_j(s,\omega),c(u)\ra ds=\int_0^{qa}


c(u_j(s,\omega))ds. \tag{2.8}


\end{equation}





{\bf 3.} В этом пункте приведем еще ряд свойств введенных отображений


$\Delta_j\in S(\Bbb R\times\Omega,\rpm(U))$


и $u_j \in S(\Bbb R\times\Omega,U)$, $j\in\Bbb N$,


отвечающих $\mu\in S(\Bbb R\times\Omega,\rpm(U))$ в


предположении, что для этого~$\mu$ уравнение (1.3) при каждом $\mm$


имеет п.\,п. по Бору (или, что равносильно, ограниченное на


$\Bbb R$) решение $x(t,\omega)=\int\limits_0^t\la \mu(s,\omega),c(u)\ra ds$.


\medskip





{\bf Лемма 3.1.} {\it


Допустим, что при каждом $\mm$\;


$\{f_j(\cdot,\omega)\}_{j=1}^{\infty}


\subset L_1^{\loc}(\Bbb R,\Bbb R^n)$, и существуют такие


строго возрастающие последовательности $\{j(l)\}_{l=1}^{\infty}$,


$\{q(l)\}_{l=1}^{\infty}\subset \Bbb N$, что при каждом $l\in\Bbb N$


\begin{equation}


\bigg|\int_0^{q(l)}f_{j(l)} (s,\omega_l) ds\bigg|\ge l. \tag{3.1}


\end{equation}


Тогда найдутся строго возрастающая последовательность


$\{h_k\}_{k=1}^{\infty}\subset (0,\infty)$ и константа


$\gamma >0$ такие, что}


\begin{equation}


\inf\limits_{N\in\Bbb N}\bigg( \sup\begin{Sb}


k,l\ge N \\ p\in \Bbb N \end{Sb}


\bigg(\sup\limits_{\omega\in\Omega}\bigg(\sup\limits_{t\in


\Bbb R}\bigg|\int_{t+h_k}^{t+h_{k+p}}f_{j(l)}(s,\omega)ds


\bigg|\bigg)\bigg)\bigg)\ge\gamma.


\tag{3.2}


\end{equation}





{\bf Доказательство.} Допустив


противное, получим, что для последовательности $h_k\doteq


k$, $k\in \Bbb N$, и константы $\gamma \doteq 1/4$


найдется такое $\wh N\in \Bbb N$, что при всех $k,l\ge \wh N$ и $p=k$


будет выполнено неравенство


$$


\sup\limits_{t\in


\Bbb R}\bigg|\int_{t+k}^{t+2k}f_{j(l)}(s,\omega_l)ds\bigg|<\frac{1}{2}.


$$


Теперь, взяв $l_{0}\ge\wh N$ таким, чтобы $q(l_{0})\ge\wh N$,


принимая во внимание (3.1), получим противоречивые неравенства


$$


l_{0}\le \bigg|\int_{0}^{q(l_{0})}f_{j(l_0)}(s,\omega_{l_0})ds\bigg|<


\frac{1}{2}.\quad\square


$$





{\bf Лемма 3.2.} {\it


Допустим, что 


$\{f_j(\cdot,\omega)\}_{j=1}^{\infty}\subset


L_1^{\loc}(\Bbb R,\Bbb R^n)$ при каждом $\mm$,


и существуют константа $\kappa>0$ и


строго возрастающие последовательности


$\{j(l)\}_{l=1}^{\infty}$,


$\{q(l)\}_{l=1}^{\infty}\subset\Bbb N$ такие, что при каждом


$l\in \Bbb N$


\begin{equation}


\bigg|\int_0^{q(l)}f_{j(l)}(s,\omega_l)ds\bigg|\ge \kappa. \tag{3.3}


\end{equation}


Тогда найдутся строго возрастающая последовательность


$\{h_k\}_{k=1}^{\infty}\subset (0,\infty)$ и константа


$\gamma$, принадлежащая $(0,\kappa)$, для которых будет выполнено


неравенство} (3.2).


\medskip





Для доказательства леммы, предположив противное,


возьмем последовательность $h_k\doteq k$, $k\in\Bbb N$,


константу $\gamma\doteq \kappa/8$ и


так же, как при доказательстве предыдущей леммы, придем к


неравенству, противоречащему (3.3).





Aналогично доказываются следующие два утверждения.


\medskip





{\bf Лемма 3.3.} {\it


Допустим, что 


$\{f_j(\cdot,\omega)\}_{j=1}^{\infty}\subset


L_1^{\loc}(\Bbb R,\Bbb R^n)$ при каждом $\mm$, и существуют такие


строго возрастающая $\{j(l)\}_{l=1}^{\infty}\subset


\Bbb N$ и строго убывающая $\{q(l)\}_{l=1}^{\infty}\subset\Bbb Z_-$


последовательности, что при каждом $l\in N$ выполнено


неравенство $(3.1)$. Тогда найдутся строго убывающая


последовательность $\{h_k\}_{k=1}^{\infty}\subset


(-\infty,0)$ и константа $\gamma>0$, для которых


будет выполняться неравенство} (3.2).


\medskip





{\bf Лемма 3.4.} {\it


Допустим, что 


$\{f_j(\cdot,\omega)\}_{j=1}^{\infty}\subset


L_1^{\loc}(\Bbb R,\Bbb R^n)$ при каждом $\mm$,


и существуют константа $\kappa>0$, а также


строго возрастающая $\{j(l)\}_{l=1}^{\infty}\subset


\Bbb N$ и строго убывающая


$\{q(l)\}_{l=1}^{\infty}\subset\Bbb Z_-$


последовательности такие, что при каждом $l\in N$ выполнено


неравенство $(3.3)$. Тогда найдутся строго убывающая


последовательность $\{h_k\}_{k=1}^{\infty}\subset


(-\infty,0)$ и константа $\gamma\in (0,\kappa)$, для которых


будет выполнено неравенство} (3.2).


\medskip





Теперь при каждом $j\in \Bbb N$ рассмотрим отображение


\begin{equation}


\tm\mapsto f_j\tm\doteq \la \Delta_j\tm,c(u)\ra\in\Bbb R^n,


\quad \tm\in \Bbb R\times\Omega,


\tag{3.4}


\end{equation}


принадлежащее (см. определение 1.2) $S(\Bbb R\times\Omega,\Bbb R^n)$.


\medskip





{\bf Лемма 3.5.} {\it


Пусть 


$\sup\limits_{t\in\Bbb R}|x\tm|<\infty$ при каждом $\mm$. Тогда}


$$


\underset{j\to\infty}{\varlimsup}


\bigg(\sup\limits_{\mm}\bigg(\sup\limits_{t\in\Bbb R}


\bigg|\int_0^t f_j(s,\omega)ds\bigg|\bigg)\bigg)<\infty.


$$





{\bf Доказательство.} Так как


(см. определение мер $\Delta_j$ и (3.4))


$\sup\limits_{\tm\in\Bbb R\times\Omega}


|f_j\tm|\le \frak c$, $j\in\Bbb N$, где (всюду далее)


$\frak c\doteq \max\limits_{u\in U}|c(u)|$, то для


доказательства леммы 3.5 достаточно показать, что


$$


\underset{j\to\infty}{\varlimsup}\bigg(


\sup\limits_{\mm}\bigg(\sup\limits_{q\in\Bbb Z}\bigg|\int_0^q


f_j(s,\omega)ds\bigg|\bigg)\bigg)<\infty .


$$


Если допустить противное, то найдутся последовательности


$\{\omega_l\}_{l=1}^{\infty}\subset\Omega$,


$\{q(l)\}_{l=1}^{\infty}\subset \Bbb Z$,


а также строго возрастающая последовательность


$\{j(l)\}_{l=1}^{\infty}$ такие, что при каждом


$l\in\Bbb N$ для функций $f_{j(l)}$, определенных равенством (3.4),


будет выполнено неравенство (3.1), которое в силу условия леммы 3.5 и


предельного соотношения (2.5) невозможно, если


$\sup\limits_{l\in \Bbb N}|q(l)|<\infty$,


т.\,е. при  сделанном предположении необходимо


$\lim\limits_{l\to\infty}|q(l)|=\infty$. Поэтому возможны


случаи, когда из последовательности $\{q(l)\}_{l=1}^{\infty}$ можно выделить


строго возрастающую подпоследовательность, принадлежащую $\Bbb N$, либо


строго убывающую подпоследовательность


из $\Bbb Z_-$. Рассмотрим первый возможный


случай, и для простоты будем считать,


что сама последовательность $\{q(l)\}_{l=1}^{\infty}\subset


\Bbb N$ является строго возрастающей. Тогда по лемме 3.1


найдутся строго возрастающая последовательность


$\{h_k\}_{n=1}^{\infty}\subset(0,\infty)$ и константа


$\gamma>0$ такие, что будет выполнено неравенство (3.2). Из


него вытекает, что для каждого $N\in \Bbb N$


найдутся такие $k(N), l(N)\ge N$ и $p(N)\in\Bbb N$, при которых


\begin{equation}


\sup_{\mm}F_{N}(\omega)\ge \frac{\gamma}{2},\quad


F_{N}(\omega)\doteq


\sup\limits_{t\in\Bbb R}\bigg|\int_{t+a(N)}^{t+b(N)}


f_{\frak j(N)}(s,\omega)ds\bigg|,\tag{3.5}


\end{equation}


где


$a(N)\doteq h_{k(N)}$, $b(N)\doteq h_{k(N)+p(N)}$,


${\frak j(N)}\doteq j(l(N))$,


$N\in \Bbb N$. Из (3.5) получаем, что для всякого $N\in \Bbb N$\;


$b(N)-a(N)\ge \frak l\doteq\frac{\gamma}{2\frak c}$ и, кроме того,


найдется такая последовательность


$\{\omega(N)\}_{N=1}^{\infty}\subset\Omega$, что при каждом $N\in\Bbb N$\;


$F_{N}(\omega(N))>\gamma/4$. Далее, т.\,к. $(\Omega,\rho)$ --- компактное


метрическое пространство, то без ограничения общности можно считать,


что $\omega(N)\to\wh\omega\in\Omega$ при $N\to\infty$. Теперь из


неравенства $F_{N}(\omega(N))>\gamma/4$ (см. (3.4), (3.5)),


$N\in\Bbb N$, получаем, что при всех $N\in\Bbb N$


\begin{equation}


\frac{\gamma}{4}\le \sup\limits_{\mm}


\bigg(\sup\limits_{t\in\Bbb R}\int_t^{t+\frak l}


|\langle\eta_{\frak j(N)}^{(1)}(s,\omega),c(u)\rangle| ds\bigg)+


I_{\frak l}(N)+ \sup\limits_{t\in\Bbb R}\bigg|\int_{t+a(N)}^{t+b(N)}\la


\mu(s,\wh\omega),c(u)\ra ds\bigg|,\tag{3.6}


\end{equation}


где


$$


I_{\frak l}(N)\doteq\sup\limits_{t\in \Bbb R}\int_{t}^{t+\frak l}


|\la\mu(s,\omega(N))-\mu(s,\wh\omega),c(u)\ra| ds.


$$


Покажем, что $\lim\limits_{N\to\infty}I_1(N)=0$. В самом деле, т.\,к.


$\mu\in S(\Bbb R\times\Omega,\rpm(U))$, то отображение (см. определение 1.2)


$\tm\to\gf\tm\doteq\la\mu\tm,c(u)\ra$


принадлежит $S(\Bbb R\times\Omega,\Bbb R^n)$ и,


следовательно


(см. определение 1.1 и (1.2)), для заданного $\varepsilon>0$


найдется такое $\gamma_0>0$,


что $\frak d_{\gamma_0}[\gf,\Omega]<\varepsilon$.


Так как $\rho(\omega(N),\wh\omega)\to 0$ при $N\to\infty$, то найдется такое


$N_{0}$, что при всех $N\ge N_0$\;


$\rho(\omega(N),\wh\omega)\le\gamma_{0}$, и поэтому


$I_1(N)\le\frak d_{\gamma_0}[\gf,\Omega]<\varepsilon$, т.\,е.


$\lim\limits_{N\to\infty}I_1(N)=0$. Отсюда,


используя топологическую эквивалентность $d_l$-расстояний ([1], с.\,198),


получаем $\lim\limits_{N\to\infty}I_{\frak l}(N)=0$. Тогда 


в силу (2.5) из (3.6) вытекает, что найдется такое $\wh N\in \Bbb N$, что


при всех $N\ge\wh N$\;


$\sup\limits_{t\in\Bbb R}|x(t+b(N),\wh\omega)-x(t+a(N),\wh\omega)|


\ge\gamma/8$.


Последнее означает, что множество п.\,п.


по Бору функций $\{x(\cdot+\gh_n,\wh\omega)\}_{n=1}^{\infty}$, где


$\gh_1\doteq a(\wh N)$, $\gh_2\doteq b(\wh N)$,


$\gh_3\doteq a(\wh N+1)$, $\gh_4\doteq b(\wh N+1),\dotsc$,


не имеет конечной $\gamma/8$-сети, что в силу


теоремы Бохнера [1] невозможно.


Точно так же, используя


лемму 3.3, придем к противоречию в случае, когда из


последовательности $\{q(l)\}_{l=1}^{\infty}$ можно выделить


строго убывающую подпоследовательность, содержащуюся в $\Bbb Z_-$.





Из леммы 3.5 и теоремы об интеграле от п.\,п. по Степанову функций


([1], с\,.206) с учетом неравенства (см. (3.4)) $|f_j\tm|\le\frak c$,


$j\in\Bbb N$, вытекает


\medskip





{\bf Следствие 3.1.} Пусть при каждом $\mm$\;


$\sup\limits_{t\in\Bbb R}|x\tm|<\infty$. Тогда существует


строго возрастающая последовательность


$\{j(l)\}_{l=1}^{\infty}\subset \Bbb N$ такая, что при любых


$l\in \Bbb N$ и $\mm$ уравнение $\Dot x=\la \Delta_{j(l)}\tm,c(u)\ra$


имеет п.\,п. по Бору решение


\begin{equation}


\frak x_{j(l)}\tm=\int_0^t \la \Delta_{j(l)}(s,\omega),c(u)\ra ds,\quad


t\in \Bbb R,\tag{3.7}


\end{equation}


причем множество $\{\frak x_{j(l)}(\cdot,\omega),\ l\in\Bbb N,\ \mm\}\subset


B(\Bbb R,\Bbb R^n)$ ограничено.


\medskip





Из следствия 3.1 и равенства (2.8) получаем


\medskip





{\bf Следствие 3.2.} Пусть при каждом $\mm$\;


$\sup\limits_{t\in\Bbb R}|x\tm|<\infty$ и


$\{j(l)\}_{l=1}^{\infty}$ --- последовательность из следствия


3.1. Тогда при любых $l\in\Bbb N$ и $\mm$ уравнение


$\Dot x =c(u_{j(l)}\tm)$ имеет п.\,п. по Бору решение


\begin{equation}


x_{j(l)}\tm=\int_0^t c(u_{j(l)}(s,\omega))ds, \tag{3.8}


\end{equation}


и при этом множество


$\{x_{j(l)}(\cdot,\omega), l\in \Bbb N, \mm\}\subset B(\Bbb R,\Bbb R^n)$


ограничено.


\medskip





{\bf 4.} При каждом $l\in\Bbb N$ рассмотрим (см. следствия 3.1, 3.2)


п.\,п. по Бору функции $\Delta x_{j(l)}(\cdot,\omega)\doteq


x(\cdot,\omega)-\frak x_{j(l)}(\cdot,\omega)$,


$\Delta \frak x_{j(l)}(\cdot,\omega)\doteq


x_{j(l)}(\cdot,\omega)-\frak x_{j(l)}(\cdot,\omega)$,


$\mm$.


\medskip





{\bf Лемма 4.1.} {\it Пусть при каждом $\mm$\;


$\sup\limits_{t\in\Bbb R}|x\tm|<\infty$. Тогда}


$$


\underset{l\to\infty}{\varlimsup}


(\sup\limits_{\mm}(\sup\limits_{q\in\Bbb Z}|


\Delta x_{j(l)}(q,\omega)|))=0.


$$





{\bf Доказательство.} Допустим


противное. Тогда найдутся такие константа $\kappa>0$,


последовательности


$\{\omega_l\}_{l=1}^{\infty}\subset\Omega$,


$\{q(l)\}_{l=1}^{\infty}\subset\Bbb Z$,


строго возрастающая подпоследовательность $\{\gj


(l)\}_{l=1}^{\infty}\subset\{j(l)\}_{l=1}^{\infty}$,


что при каждом $l\in\Bbb N$ для функции (см. (3.7))


\begin{equation}


\tm\to


f_{\gj(l)}\tm\doteq \la


\eta_{\gj(l)}^{(1)}\tm,c(u)\ra\in \Bbb R^n,\quad


\tm\in\Bbb R\times\Omega, \tag{4.1}


\end{equation}


будет выполнено неравенство (3.3) при $j(l)\doteq \gj(l)$,


которое в силу предельного соотношения (2.5) невозможно, если


$\sup\limits_{l\in\Bbb N}|q(l)|<\infty$,


т.\,е. при сделанном предположении необходимо


$\lim\limits_{l\to\infty}|q(l)|=\infty$. Поэтому возможны


случаи, когда из этой последовательности можно выделить


строго возрастающую подпоследовательность, принадежащую $\Bbb N$, либо


строго убывающую подпоследовательность из $\Bbb Z_-$.


В первом случае считаем,


что сама последовательность


$\{q(l)\}_{l=1}^{\infty}\subset\Bbb N$ является строго


возрастающей. По лемме 3.2 найдутся строго


возрастающая последовательность


$\{h_k\}_{k=1}^{\infty}\subset(0,\infty)$ и константа


$\gamma\in (0,\kappa)$ такие, что для функции


$f_{{\frak j}(l)}$,


определенной равенством (4.1), при $j(l)=\gj(l)$ будет


выполнено неравенство (3.2).


Из него вытекает, что для этой


функции и для каждого $N\in\Bbb N$ найдутся такие $k(N),l(N)\ge N$ и


$p(N)\in\Bbb N$, что будет иметь место неравенство (3.5) при


${\frak j}(N)={\frak i}(N)\dot={\frak j}(l(N))$.


Отсюда для всякого $N\in\Bbb N$\;


$b(N)-a(N)\ge\frac{\gamma}{4\frak c}$ и существует такая


последовательность точек $\{\omega_N\}_{N=1}^{\infty}\subset\Omega$, что


$F_{N}(\omega_N)\ge \gamma/4$. Следовательно, выполняется неравенство


$$


\frac{\gamma}{4}\le\sup\limits_{\mm}\bigg(\sup\limits_{t\in\Bbb R}


\int_{t}^{t+\frak l}


|\la\eta_{{\frak i}(N)}^{(1)}(s,\omega_N),c(u)\ra| ds\bigg),


$$


которое не совместно с (2.5)


при $\frak l\doteq\frac{\gamma}{4\frak c}$. Точно так же, используя


лемму 3.4, придем к противоречию в


случае, когда из последовательности


$\{q(l)\}_{l=1}^{\infty}$ можно выделить строго убывающую


подпоследовательность, содержащуюся в $\Bbb Z_-$.


\medskip





{\bf Доказательство теоремы 1.2.} В


силу леммы 4.1 из последовательности $\{j(l)\}_{l=1}^{\infty}$


можно выделить такую подпоследовательность $\{\gj(i)\}_{i=1}^{\infty}$,


что (см. следствие 3.2) при каждом $i\in\Bbb N$ и всяком $\mm$


функция $x_i(\cdot,\omega)\doteq x_{\gj(i)}(\cdot,\omega)$


--- п.\,п. по Бору решение уравнения $\Dot x=c(v_i\tm)$,


где $v_i\tm\doteq u_{j(i)}\tm$, $\tm\in\Bbb R\times\Omega$. При этом


множество функций $\{x_i(\cdot,\omega),i\in\Bbb N,\mm\}\subset


B(\Bbb R,\Bbb R^n)$ ограничено и


\begin{equation}


\lim\limits_{i\to\infty}(\sup\limits_{\mm}


(\sup\limits_{m\in\Bbb Z}|\Delta x_{\gj(i)}(m,\omega)|))=0. \tag{4.2}


\end{equation}


При любых $q\in\Bbb N$, $i\in\Bbb N$


и всяком $\mm$ (см. (2.8), (3.7) и


(3.8)) имеем соотношения


\begin{multline*}


\frac{1}{qa}\int_0^{qa}


|x\tm-x_{\gj(i)}\tm|dt\le\frac{1}{qa}


\sum\limits_{m=0}^{q-1}\int_{ma}^{(m+1)a}


(|\Delta x_{\gj(i)}\tm|+|\Delta \frak x_{\gj(i)}\tm|) dt\le \\


%


\le\sup\limits_{\mm}(\sup\limits_{m\in\Bbb Z}|\Delta


x_{\gj(l)}(m,\omega)|)+


\sup\limits_{\mm}\bigg(\sup\limits_{t\in\Bbb R}\int_t^{t+a}|\la


\eta_{\gj(i)}^{(1)}(s,\omega),c(u)\ra |ds\bigg)+ \\


%


+\sup\limits_{\mm}\bigg(\sup\limits_{m+\Bbb Z}


\bigg(\sup\limits_{t\in[ma,(m+1)a]}\bigg|\int_{ma}^t\la


\eta_{\gj(i)}^{(2)}(s,\omega), c(u)\ra ds\bigg|\bigg)\bigg),


\end{multline*}


из которых в силу предельных равенств (2.5) и (2.7) при $\frak l\doteq a$


и (4.2) получаем


$$


\lim\limits_{i\to\infty}\sup\limits_{\mm}M\{|x\tm-x_i\tm|\}\doteq


\lim\limits_{i\to\infty}\sup\limits_{\mm}M\{|x\tm-x_{\gj(i)}\tm|\}=0.


\quad\square


$$


\medskip





Отметим, что доказанная теорема 1.2 играет важную роль, во-первых, при


обосновании корректности расширения задач оптимального управления п.\,п.


движениями (напр., [8]) с уравнениями связи вида


$\Dot x=c(u(t))$, $t\in\Bbb R$, $u(\cdot)\in S(\Bbb R,U)$, а


во-вторых, при доказательстве необходимых условий оптимальности в таких


задачах.





\begin{thebibliography}{9}





\bibitem{1}


Левитан Б.М. {\em Почти периодические функции}. -- М.: Гостехиздат, 1953.


-- 396 с.


\bibitem{2}


Fink A.M. {\em Almost periodic differential equation} // Lect. Notes Math.


-- 1973. -- V.\,377. -- 336~p.


\bibitem{3}


Варга Дж. {\em Оптимальное управление дифференциальными и функциональными


уравнениями}. -- М.: Наука, 1977. -- 623~с.


\bibitem{4}


Иванов А.Г. {\em Мерозначные почти периодические функции}. -- Препринт.


ФТИ УрО АН СССР, 1990. -- 53~с. 


\bibitem{5}


Иванов А.Г. {\em Мерозначные почти периодические функции}. --


Удмуртcк. ун-т. -- Ижевск,


1991. 62 с. -- Деп. в ВИНИТИ 24.04.91, \No\,1721-B91.


\bibitem{6}


Иванов А.Г. {\em О ляпуновской почти периодической задаче} // ПММ. --


1991. -- Т.\,55. -- Bып.\,5. -- С.\,718--724.


\bibitem{7}


Иванов А.Г. {\em Об эквивалентности дифференциальных включений и


управляемых почти периодических систем} // Дифференц. уравнения. -- 1997.


-- Т.\,33. -- \No\,7. -- С.\,876--884.


\bibitem{8}


Иванов А.Г. {\em Об оптимальном управлении почти периодическими


движениями} // ПММ. -- 1992. -- Т.\,56. -- Bып.\,5. -- С.\,718--724.





\end{thebibliography}





\vskip 0.5cm





\post{\it Удмуртский государственный}{\it Поступила}


\post{\it университет}{12.03.2001}





\label{end}


\end{document}


